\scanpage{54}
% archive sheet number: 53
Algèbres étales et la différente \mspage{(3)}\hfill Tribu n\textsuperscript{o} 65

sur $A$ tel que, en posant $B'=A[x]$, $F'(x)$ soit inversible dans $B'_{\handwritten{\unsure{q|\mathfrak{q}}\cap}B'}$ et
$B'_{\handwritten{q\cap}B'}\longrightarrow B_{\handwritten{q}}$ soit bijectif. Généralisation au cas où $B$ est localisée et une
$A$-algèbre finie.

Si, de plus, $A_{\handwritten{q\cap}A}$ est normal et si $B$ est étale en $\handwritten{q}$, $B_{\handwritten{q}}$ est normal.
La prop. 4 (p. 4) du texte.

Ajouter, entre astérisques, les améliorations dépendant du main-theorem.

\uline{n\textsuperscript{o} 5}. \uline{Groupes d'automorphismes}.

La prop. 6 du texte, dans le cas noethérien non réduit (cf. Tribu n\textsuperscript{o} 64). Son cor.
(p. 7) est la prop. appliquée à $H$, de sorte qu'on le vide.

\uline{n\textsuperscript{o} 6}. \uline{Un critère de non ramification}.

Il s'agit de l'actuel n\textsuperscript{o} 4 (\S~1) amélioré comme suit. Remplacer (ici et partout)
l'idiote formule $(A^{*})^{-1}B=L$ par $KB=L$. Référer à Alg. IX, \S~2 pour la définition
du discriminant.

Dans le b) du th. 1, ajouter

(iii)\qquad $B$ est finie sur $A$ et $B/\handwritten{\unsure{m|\mathfrak{m}}}B$ est étale sur $A/\handwritten{m}$ ;

(iv)\qquad $B$ est étale finie sur $A$ ;

comme conditions équivalentes à (i) et (ii). Ça absorbe la partie intéressante de c),
ainsi que le cor. 1.

Pour démontrer le lemme (p. 10-11) on remarque que $P(x)=D(1,x,\ldots,x^{n-1})$ est
fonction polynôme de $x$, et qu'on a $P\handwritten{\not}=0$ (passer à la clôture algébrique). On a
donc $x\handwritten{\in}R$ tel que $P(x)\handwritten{\not}=0$.

Le cor. 1 est absorbé. La déf. 2 est supprimée. Le cor. 2 devient une prop. donnant
l'équivalence de : a)\quad $B$ est étale finie sur $A$ ;

b)\quad $B_{\handwritten{m}}$ est étale finie sur $A_{\handwritten{m}}$ pour tout idéal maximal $\handwritten{m}$ de $A$ ;

c)\quad $B$ est finie sur $A$ et $B/\handwritten{m}B$ est étale sur $A/\handwritten{m}$ pour tout
idéal maximal $\handwritten{m}$ de $A$.

Le cor. 3 est vidé. On donne enfin le cor. 4 (y remplacer ``non-ramifié'' par
``étale'').

On essaiera de mettre l'\uline{appendice} suivant : si $I$ est un idéal nilpotent de $A$,
il y a une équivalence de catégories entre les algèbres étales finies sur $A$ et les
algèbres étales finies sur $A/I$. Voir S. G. A.

